% ============================================================
%  SCS 3312 – Research Methods
%  D0: Topic Proposal & Research Plan
%  Group Submission
% ============================================================

\documentclass[12pt,a4paper]{report}

% ============================================================
%  PACKAGES
% ============================================================

\usepackage[left=1.5in,right=1in,top=1.5in,bottom=1in]{geometry}
\usepackage[T1]{fontenc}
\usepackage[utf8]{inputenc}
\usepackage{setspace}
\usepackage{titlesec}
\usepackage{tocloft}
\usepackage{fancyhdr}
\usepackage{hyperref}
\usepackage{bookmark}
\usepackage{parskip}
\usepackage{enumitem}
\usepackage{booktabs}
\usepackage{array}
\usepackage{makeidx}
\usepackage{natbib}

\makeindex

% ============================================================
%  HYPERREF SETTINGS
% ============================================================

\hypersetup{
    colorlinks=true,
    linkcolor=black,
    citecolor=black,
    urlcolor=blue,
    bookmarks=true,
    bookmarksopen=true,
    pdftitle={Passenger Waiting Behaviour at Thummulla Bus Stop},
    pdfauthor={SCS3312 Group}
}


% ============================================================
%  CHAPTER STYLE
% ============================================================

\titleformat{\chapter}[block]
{\normalfont\Large\bfseries}
{\chaptername\ \thechapter}{1em}{}

\titlespacing*{\chapter}{0pt}{30pt}{20pt}

% ============================================================
%  LINE SPACING
% ============================================================

\onehalfspacing

% ============================================================

\documentclass[12pt, a4paper]{report}

\usepackage[left=1.5in, right=1in, top=1.5in, bottom=1in]{geometry}
\usepackage{setspace}
\usepackage{fancyhdr}

\pagestyle{fancy}
\fancyhf{}
\renewcommand{\headrulewidth}{0pt}
\cfoot{\thepage}

\begin{document}

% ============================================================
%  TITLE PAGE
% ============================================================
\begin{titlepage}
    \thispagestyle{empty}
    \begin{center}

        \vspace*{0.8cm}

        {\Large\bfseries UNIVERSITY OF COLOMBO SCHOOL OF COMPUTING\par}
        \vspace{0.2cm}
        {\large SCS 3312 -- Research Methods\par}

        \vspace{1.4cm}

        \rule{0.78\textwidth}{0.5pt}\par
        \vspace{0.4cm}

        {\LARGE\bfseries Study of Passenger Waiting Behaviour\par}
        \vspace{0.2cm}
        {\Large at Thummulla Bus Stop\par}

        \vspace{0.4cm}
        \rule{0.78\textwidth}{0.5pt}\par

        \vspace{1.0cm}

        {\large\bfseries Topic Proposal \& Research Plan\par}
        \vspace{0.15cm}

        \vspace{0.15cm}
        {\large Group Number: 02\par}

        \vspace{0.9cm}

        {\large\bfseries Group Members\par}
        \vspace{0.35cm}

        \begin{tabular}{@{}p{7.5cm}r@{}}
            R. R. K. Jinendra      & 23000821 \\[0.2cm]
            S. S. Balahewa         & 23000155 \\[0.2cm]
            G. C. K. S. Gajanayake & 23000481 \\
        \end{tabular}

        \vfill



    \end{center}
\end{titlepage}% ============================================================
%  FRONT MATTER
% ============================================================
\pagenumbering{roman}
% No section numbering before Chapter 1
\setcounter{secnumdepth}{-1}

% ── Acknowledgement ─────────────────────────────────────────
\chapter*{Acknowledgement}
\addcontentsline{toc}{chapter}{Acknowledgement}

We would like thanks to Dr. Manjusri Wickramasighe to his guidance  helped us to define our topic proposal and research plan. We are willing to continue our research based on your feedback.
\clearpage

\chapter*{Abstract}
\addcontentsline{toc}{chapter}{Abstract}
This research proposal aims to presents passenger behavior and boarding decisions at mentioned bus stops. This will helps to optimize bus schedules at peak and off times on passenger behavior.  As a result of that we can present many solutions to existing problems such as reducing the number of waiting passengers at the peak hours.

\clearpage
\pagenumbering{roman}
\setcounter{page}{1}
\setcounter{secnumdepth}{-1}

% ============================================================
%  TABLE OF CONTENTS
% ============================================================

\tableofcontents
\clearpage

% ============================================================
%  MAIN MATTER
% ============================================================

\pagenumbering{arabic}
\setcounter{page}{1}
\setcounter{secnumdepth}{3}


% ============================================================
%  CHAPTER 1
% ============================================================

\chapter{Introduction}

This research study looks at how passengers behave while waiting at bus stops considering peak hours, weather conditions etc. This study aims to investigate passengers with boarding decisions at bus stops identifying common behaviours of the passengers.


% ============================================================

\section{Background of the Study}

Public transportation is very important for people to do day today activities in a urban city like Colombo. Many students, office workers and other public people widely uses buses everyday because it is cheap and easy to access. So the bus stops become crowded where many people wait, move and make boarding decisions.

Passenger behaviour at bus stops changes with many different situations such as crowd level, waiting time for a bus, availability of buses, environmental conditions and urgency to go to the destinations. Some people wait for a long time for the bus they want while others take any available bus disregarding the crowdness.

Thummulla bus stop is one of the busy bus stops in Colombo. A large number of passengers use this bus stop daily, specially during morning and  evening peak hours. 

Studying these different behaviours helps to understand how people react and adapt to crowded public transport environments.

% ============================================================

\section{Research Problem}

At busy bus stops, passengers behave differnetly according to the situation. Many passengers have to face to the problems like long waiting time, irregular bus arrivals, overcrowded buses, delays and insufficient number of buses during peak hours. So these issues often affect changing passenger behaviour and making travel decisions at bus stops.

Some passengers wait for a long time for a specific bus because some people may have their season tickets. Even if the bus is overcrowded, they still get on the bus, while others choose alternative bus or differnet transport methods due to overcrowding or delays.

These passenger behaviour directly affects the efficiency of bus transportation and management of the bus schedules.
Currently bus schedules are not often adjusted according to actual passenger behaviour or crowd conditions. As a result of this overcrowding occurs at major bus stops.

Therefore, this research aims to study passenger behaviour at mentioned bus stops and identify how bus time schedules can be optimized based on passenger behaviour and peak hour travel patterns to reduce crowd at busy bus stops. 

% ============================================================

\section{Justification of the Study}

Passenger behaviour at bus stops is important because many  people use buses everyday travelling to schools, universities, work and other places. At busy bus stops people often face issues like long waiting times, overcrowded buses and delayed bus arrivals.

Some people become impatient, some changes buses and others choose different transport methods. 

Thummulla bus stop is one of busy bus stop in Colombo and experience heavy passenger traffic specially during morning and evening hours. So passengers have to face overcrowding and delays often.

Studying passenger behaviour is important  so that it helps to understand when bus stops become crowded, how passengers react to bus delays and overcrowding.

The findings of the study can help improve bus time schedule based on peak hour demand. This may help improve the overall public transport experience for passengers.  

% ============================================================

\section{Research Topic}

\begin{center}
\textit{
“Study of Passenger Waiting Behaviour at Thummulla Bus Stop”
}
\end{center}

This study focuses on understanding passenger behaviour  mainly at Thummulla Bus Stop and the nearst two other bus stops. Passenger patterns can be used to improve bus time schedules via reducing waiting time of the passengers.

%


% ============================================================

\section{Preliminary Research Questions}

This study tries to get answers for the following research questions:

\begin{enumerate}[noitemsep]

\item What problems passengers commonly face during peak hours?

\item What are common passenger behaviours at mentioned bus stops during waiting periods?

\item How do waiting times affect passenger decisions and actions?

\item Why do some passengers reject overcrowded buses?

\item How do passengers adapt to delays and overcrowding during peak hours?

\item During which time mentioned bus stops become most crowded?

\end{enumerate}

% ============================================================

\section{Scope of the Study}

This study is limited to observing passenger waiting behaviour and boarding decisions at
mentioned Bus Stops.

The research focuses on:

\begin{itemize}[noitemsep]

\item Passenger waiting time.

\item Boarding and rejection behaviour.

\item Crowd density of both field and buses.

\item Passenger movement within waiting areas

\item Peak hour and off peak hour behavioural differences.

\end{itemize}

The study does not investigate government transport policy.

% ============================================================
%  CHAPTER 2
% ============================================================

\chapter{Field Study Description}

This chapter explains the selected field site and describes the environmental conditions,
passenger activities, and behavioural patterns observed during preliminary field visits.

% ============================================================

\section{Selected Field Site}

We selected the research sites based on 3 main locations: 
\begin{itemize}[noitemsep]
\item Thummulla Bus Stop
\item University Bus Stop (Near Art Faculty)
\item University Bus Stop (Near Science Faculty).
\end{itemize}



% ============================================================

\section{Description of the Field Site}

The selected locations are very busy bus stops among university students, office workers, and the general public of Colombo

The location experiences:

\begin{itemize}[noitemsep]

\item High passenger density

\item Buses arrive frequently because these bus stops are located on highly crowded routes

\item Overcrowding in peak hours

\item Different types of people arrive at these bus stops


\item Waiting Facilities are in these bus stops


\item Traffic congestion


\end{itemize}

Based on these conditions, we can observe passenger behaviors related to waiting, movement, and boarding decisions.

% ============================================================

\section{Peak,  Off-Peak and environmental Conditions}

We normally identify that a passenger’s behaviour varies depending on the time of the day and some environmental conditions, such as the weather.
% ============================================================

\subsection{Peak Hours}

During morning and evening peak hours, a lot of people come to the bus stop, and we can observe the following conditions,

\begin{itemize}[noitemsep]

\item Passenger density increases

\item Bus arrival frequency increases

\item How passengers decide which bus to board

\item Passengers may display impatience or aggressive boarding behaviour.

\end{itemize}

% ============================================================

\subsection{Off-Peak Hours}

In this time period, passenger arrivals can decrease during the day and at night.

\begin{itemize}[noitemsep]

\item Passenger flow decreases

\item Waiting times vary 

\item Boarding behaviour of off-peak passengers


\end{itemize}


% ============================================================


\subsection{Various Weather conditions}
We can observe passengers’ behavior varies under different weather conditions, such as Rainy days and sunny days. 

\section{Observed Passenger Activities}

Preliminary observations identified several common passenger activities:

\begin{itemize}[noitemsep]

\item Waiting for specific bus routes;

\item Rejecting overcrowded buses;

\item Queue changing and movement within waiting areas;

\item Standing near the roadside before buses arrive;

\item Using mobile phones while waiting;

\item Interacting with other passengers;

\item Aggressive boarding during crowded situations.

\end{itemize}


% ============================================================
%  CHAPTER 3
% ============================================================

\chapter{Research Methodology}

In this section explains how research procedures use for collecting qualitative and quantitative data.

% ============================================================

\section{Research Approach}

This study used a mixed method with combining qualitative and quantitative data collection methods.

Qualitative methods will be used to understand passenger behaviour and experiences,
whereas quantitative methods will used to measure numerical data for the waiting
time and crowd conditions.

% ============================================================

\section{Qualitative Data Collection Methods}

The qualitative focuses on identifying passenger behaviours.

% ============================================================

\subsection{Observation}

Observation are  used to study:

\begin{itemize}[noitemsep]

\item Bus Frequency
\item Field crowd.
\item Seat Availabilty of the bus.
\item Boarding desitions on buses.
\item Bus Types
\item Is bus overcrowding
\item passenger waiting time.

\end{itemize}

% ============================================================

\subsection{Informal Interviews}

Short interviews are used with selected passengers.

Interview topics include:

\begin{itemize}[noitemsep]

\item What are the reasons for waiting for specific buses ?.
\item Passenger opinions regarding overcrowding on the bus stop.
\item Experiences during long waiting periods.
\item passenger waiting time.

\end{itemize}

\subsection{Survey Data Collection}

Questionnaires will be give to selected passengers.

questions will examine:

\begin{itemize}[noitemsep]

\item Is the passenger worker, student or someone else.
\item Priority of to go to the destinations.
\item Why not selecting any other bus insted of relevent one (like 138 passenger can board into 122 or 125).


\end{itemize}

% ============================================================

\section{Quantitative Data Collection Methods}

The quantitative focuses on measurable numerical data.

% ============================================================

\subsection{Passenger Counting}

Passenger counting will record.

\begin{itemize}[noitemsep]

\item Number of waiting passengers on mentioned bus stops.
\item Number of passengers get into the bus.
\item Number of bus rejection by passengers.

\end{itemize}

% ============================================================

\subsection{Waiting Time Measurement}

Waiting-time measurements will determine:

\begin{itemize}[noitemsep]

\item Average waiting time;

\item Maximum waiting time;

\item Waiting time differences during peak and off-peak periods.

\end{itemize}




% ============================================================

\section{Variables to be Measured}

The following variables will be measured:

\begin{itemize}[noitemsep]

\item Passenger waiting time.
\item Crowd density
\item Passenger count
\item How many Passengers are boarded into bus.
\item Number of rejected buses
\item Bus arrival intervals.
\item Bus type.
\item Bus seating availability.
\item Traffic condition.
\item Age group of the passenger.
\item Waiting time for the bus.
\item Gender.
\item Bus frequency.
\item Bus route(Number).
\item How many buses are missed.
\item Reasons for rejecting buses.
\item Even bus is crowed how many passengers get into bus




\end{itemize}

% ============================================================

\section{Feasibility of the Study}

This research is feasible because all three bus stops are publicly accessible and near the university and also all locations are passenger active all consideration times. 
This study will be done using simple and practical methods like Observation, passenger counting, interviews, and waiting-time measurements. 

The selected time slots have sufficient passenger volume at all day. Therefore it gives meaningfull data sample.

Therefore this study is suitable for mixed methods investigation.


\section{Ethical Considerations}

The research is conduct to meet ethical standards through all stages of the data collection. Not collect any personal informations such as name, contact details or any other identifications. Thus, passenger privacy will be respected.


% ============================================================
%  REFERENCES
% ============================================================

\clearpage
\renewcommand{\bibname}{References}
\addcontentsline{toc}{chapter}{References}

\begin{thebibliography}{9}

\bibitem{transport2023}
Ministry of Transport and Highways Sri Lanka. (2023). 
\textit{National Transport Policy}. Available at: 
\url{https://www.transport.gov.lk/web/images/downloads/F-CoMTrans_Main_S.pdf}
(Accessed: 20 May 2026).

\bibitem{surveydoc}
National Transport Commission (2024). 
\textit{Passenger Transport Related Document}. Available at: 
\url{https://drive.google.com/file/d/1IiPWAT_psSz5t9elF5lUlvERJNI2pgeL/view?usp=sharing}
(Accessed: 20 May 2026).

\end{thebibliography}



% ────────────────────────────────────────────────────────────
%  INDEX  (no section numbering)
% ────────────────────────────────────────────────────────────
\chapter*{Index}
\addcontentsline{toc}{chapter}{Index}

\begin{table}[h!]
\centering
\begin{tabular}{p{7cm} p{5cm}}
\toprule
\textbf{Term} & \textbf{Page(s)} \\
\midrule
Boarding behaviour          & 3, 7, 9, 11 \\
Bus rejection               & 3, 7, 9, 11, 13 \\
Bus stop                    & 4, 5, 6, 8 \\
Colombo                     & 4, 5, 8 \\
Crowd density               & 2, 3, 7, 10, 13 \\
Data collection             & 9, 10, 11, 12 \\
Ethical considerations      & 14 \\
Field notes                 & 9, 10 \\
Field sites                 & 5, 9 \\
Fort Main Bus Stop          & 5, 6, 15 \\
Informal interviews         & 9, 10, 12 \\
Mixed-method approach       & 8, 14, 15 \\
Observation                 & 9, 10, 12 \\
Off-peak hours              & 3, 6, 7, 8, 11 \\
Passenger counting          & 10, 11, 15 \\
Passenger safety            & 2, 3, 13, 14 \\
Peak hours                  & 3, 5, 6, 8, 11, 13 \\
Queue behaviour             & 2, 3, 9, 16 \\
Reid Avenue Bus Stop        & 5, 6, 15 \\
Research objectives         & 4 \\
Research questions          & 4 \\
Survey                      & 10, 11, 12, 15 \\
Town Hall Bus Stop          & 5, 6, 15 \\
Waiting time                & 2, 3, 7, 10, 11, 13, 15 \\
\bottomrule
\end{tabular}
\end{table}
\end{document}


